\documentclass[bib]{statapress}
\usepackage[crop,newcenter,frame]{pagedims}
\usepackage{sj}
\usepackage{epsfig}
\usepackage{graphicx}
\usepackage{stata}
\usepackage{amsmath}
\usepackage{threeparttable}
\usepackage{tabularx}
\sjsetissue{$vv$}{$ii$}{$mm$}{$yyyy$}
\graphicspath{{replication/figures/}}
\emergencystretch=3em

\newcommand{\gracedr}{\texttt{gracedr}}

\begin{document}

\inserttype[st0001]{article}
\author{Ma and Hardin}{%
  Grace Huiyi Ma\\Dutch Fork High School STEM Magnet Program\\Irmo, SC\\
  \and
  James W. Hardin\\Department of Epidemiology and Biostatistics\\
  University of South Carolina\\Columbia, SC\\jhardin@sc.edu
}
\title[gracedr: Randomization with exact group sizes]{%
  gracedr: Modified minimal sufficient balance randomization with
  exact group sizes, an audit log, and encrypted datasets in Stata}
\maketitle

\begin{abstract}
Minimal sufficient balance randomization balances baseline covariates while
keeping assignments unpredictable, but it does not guarantee prescribed group
sizes, and the usual remedy of capping the arm sizes forces long runs of
predictable assignments. We describe \gracedr, a Stata package that implements
a modified procedure. It adds an exact-allocation run-in, forced allocation at
enrollment totals chosen in advance and at the final sample size, any number of
arms with unequal allocation ratios, and a base probability computed from the
places left before the next forced total. The package stores each study in a
folder, records every assignment and override in an append-only audit log,
replays any past assignment, and simulates operating characteristics. Because
Stata has no built-in encrypted file format, we also describe \texttt{encsave}
and \texttt{encuse}, which save and load datasets as password-protected
encrypted zip files through Stata's Python integration, and encrypted studies built on them. We
simulated 1,000 trials in each of four settings (200 or 500 participants, with
and without drift in the enrolled population) and compared the procedure with
simple randomization, permuted blocks, minimization, and unmodified minimal
sufficient balance. The modified procedure met the group sizes at the end of the
run-in and at the end of enrollment in every trial, whereas unmodified minimal
sufficient balance missed the final sizes by a mean of 5.8 to 10.2 participants.
The mean largest standardized difference at the end of enrollment of the
modified procedure (0.093 to 0.162) was lower than for simple randomization and
permuted blocks and higher than for unmodified minimal sufficient balance and
for minimization, and its expected correct-guess probability was 58.5\% to
59.3\%, compared with 70.8\% to 70.9\% for permuted blocks of size four and
61.7\% to 62.1\% for minimization.
\keywords{\inserttag, \gracedr, randomize, \texttt{encsave}, \texttt{encuse}, covariate-adaptive
randomization, minimal sufficient balance, dynamic randomization, allocation
randomness, encryption}
\end{abstract}

\section{Introduction}

In many clinical trials participants become eligible one at a time and must be
assigned to treatment immediately. The randomization procedure must then
balance two goals that pull against each other. Treatment assignments should be
unpredictable, so that investigators cannot steer participants toward an arm,
and the arms should be comparable in size and in prognostically important
baseline covariates, so that chance imbalance does not cloud the treatment
comparison \citep{rosenberger2016,berger2005,altman1990}. Restricted designs
such as permuted blocks and biased coins control group sizes
\citep{efron1971,wei1978,soares1983,berger2003}, and covariate-adaptive designs
such as minimization control covariate margins
\citep{taves1974,pocock1975,scott2002}, but each restriction makes a larger
share of assignments predictable \citep{blackwell1957,zhao2012compare}.

Minimal sufficient balance (MSB) randomization \citep{zhao2015msb} resolves
much of this tension by using a biased coin only when a controlled covariate is
imbalanced beyond a prespecified limit \emph{and} the current participant's
assignment can reduce that imbalance. Covariates are used on their original
scale, so continuous covariates need not be categorized. Later work examined
the statistical properties, bias, and power of MSB
\citep{lauzon2020msb,lauzon2022msb} and proposed a common-scale variant
\citep{johns2022csmsb}.

MSB as published, however, does not guarantee the prescribed group sizes. Most
protocols prescribe the number of participants in each arm, and the charters of
data and safety monitoring boards may expect the arms to be close to their
prescribed sizes at each review. The simplest remedy, capping each arm at its
prescribed size, creates the problem MSB was designed to avoid: once one arm
fills, every remaining participant is forced into the other arm, and those
forced assignments are perfectly predictable. When the enrolled population
changes over time, the forced tail also turns that time trend into covariate
imbalance. \citet{zhao2015msb} noted that a balanced allocation could be
enforced for an initial portion of the trial with the random allocation rule,
but they did not formalize the idea.

This article describes the \gracedr\ package for Stata 19, which implements a
modified MSB procedure built around these operational requirements. Relative
to the original algorithm it adds an exact-allocation run-in whose size follows
a stated rule, forced allocation at enrollment totals chosen in advance and at
the final sample size, support for any number of arms and for unequal
allocation ratios, and a base probability computed from the places left before
the next forced total. A group-size vote and a big stick bound
\citep{soares1983} complete the procedure. Around the algorithm, the package
keeps the study in a folder with an append-only audit log, replays any past
assignment, and simulates operating characteristics against four comparators.

Live randomization needs a record that the investigators cannot change
unnoticed and that unauthorized people cannot read, because the allocation
sequence must stay concealed. Stata has no encrypted file format. We therefore
also built \texttt{encsave} and \texttt{encuse}, two commands that are
independent of the randomization commands, and an option that stores a whole
study as encrypted files. Section~\ref{sec:encrypt} describes their
development.

Section~\ref{sec:algorithm} states the algorithm, section~\ref{sec:commands}
describes the commands with a worked example, section~\ref{sec:encrypt}
describes the encrypted datasets, section~\ref{sec:sim} reports a simulation
study, and section~\ref{sec:conclusions} discusses what the simulations do and
do not show.

\section{The algorithm}
\label{sec:algorithm}

\subsection{Notation}

Consider a trial with $k \ge 2$ arms and target sample size $N$. The allocation
ratio is a vector of positive numbers, and the target proportions are
$r_a = \text{ratio}_a / \sum_b \text{ratio}_b$ for arms $a = 1, \ldots, k$.
Before participant $i$ is randomized, $n_a(i)$ is the number of participants
already assigned to arm $a$, so that $\sum_a n_a(i) = i - 1$, and
\[
D_a(i) = n_a(i) - r_a (i - 1)
\]
is the deviation of arm $a$ from its target count. Every allocation rule in the
package returns a probability vector $\phi(i) = (\phi_1(i), \ldots,
\phi_k(i))$, together with a label for the step of the algorithm that produced
it. The participant is assigned by one draw from $\phi(i)$.

\subsection{Minimal sufficient balance votes}
\label{sec:votes}

Each balance variable $v$ has a weight $w_v > 0$ that we think of
as a type of importance weight. That is, these weights reflect the 
proportion of that variable's contribution to calculating overall
imbalance.  The weights are rescaled to
sum to one, and equal weights are the default. Each variable casts a vote for
one arm or no vote. Votes are collected only when the imbalance between the arms
among the participants already randomized is large enough, as measured by a
$p$-value below a user-specified limit $p^\ast$ (default 0.30), and when the current
participant's value allows specific group assignment to reduce that imbalance
\citep{zhao2015msb}. A participant whose value of a variable is missing casts
no vote for that variable.

\emph{Categorical variables.} Let $T$ be the table of previously randomized participants by
level and arm. Levels with no previous participants are dropped. No vote is
cast if fewer than two levels remain, if any arm has no participants in the
table, or if the current participant's level is not in the table. Otherwise the
Pearson statistic $\chi^2 = \sum_{j,a} (T_{ja} - E_{ja})^2 / E_{ja}$, with
$E_{ja}$ the product of the row and column totals divided by the table total,
is compared with a $\chi^2$ distribution on $(L - 1)(k - 1)$ degrees of
freedom, where $L$ is the number of levels in the table. If the $p$-value is
below $p^\ast$, the vote goes to the arm with the lowest share
$T_{ja} / \sum_j T_{ja}$ in the current participant's level $j$. A tie for the
lowest share gives no vote.

\emph{Continuous variables.} No vote is cast unless every arm has at least two
previous participants and a positive variance. With two arms, the statistic is
Welch's two-sample $t$ statistic with the Satterthwaite degrees of freedom
\citep{welch1947,satterthwaite1946}. With $k > 2$ arms it is Welch's one-way
analysis of variance from the arm sizes $n_a$, means $m_a$, and variances
$s_a^2$ \citep{welch1951}. If the $p$-value is below $p^\ast$, a participant
whose value is above every arm mean votes for the arm with the smallest mean,
and a participant whose value is below every arm mean votes for the arm with the
largest mean. A participant whose value lies between the arm means casts no
vote, because either assignment would change the imbalance very little. A tie
for the smallest or largest mean gives no vote.

\emph{Group size.} After the run-in, group size is one more voter. If the arm
with the most negative deviation has $D_a(i) \le -1$, and no other arm has the
same deviation, it receives a vote of weight 1.

\emph{Tally.} The score of arm $a$ is the sum of the weights of the votes it
received. The winner $\hat a$ is the arm with the single highest score. There is
no winner if no votes were cast or if two arms tie for the highest score.

\subsection{The biased coin}

MSB assigns to the favored arm with probability $\xi$ (default 0.65). To apply
the same coin to a base probability vector $b$ that need not be uniform, the
package tilts the favored arm on the odds scale. With $\omega = \xi / (1 - \xi)$,
\begin{equation}
\phi_{\hat a} = \frac{\omega b_{\hat a}}{1 - b_{\hat a} + \omega b_{\hat a}},
\qquad
\phi_a = b_a \, \frac{1 - \phi_{\hat a}}{1 - b_{\hat a}} \quad (a \ne \hat a).
\label{eq:tilt}
\end{equation}
The other arms keep their relative probabilities. With two arms and
$b = (1/2, 1/2)$, the favored arm receives exactly $\xi$, which is the rule of
\citet{zhao2015msb}. If there is no winner, or if $b_{\hat a}$ is 0 or 1, then
$\phi = b$, so a forced assignment is never overridden by a vote.

\subsection{Checkpoints and the base probability}

A \emph{checkpoint} is an enrollment total at which the arm counts must equal
their target counts. The checkpoints $c_1 < \cdots < c_J = N$ are the end of the
run-in (when the run-in is not empty), any totals the user names in advance with
the \texttt{forcedat()} option, and the final sample size $N$. The target count
of arm $a$ at checkpoint $c$ is $t_a(c)$, the largest-remainder rounding of
$c \, r_a$: each arm receives $\lfloor c r_a \rfloor$, and the remaining units
go one at a time to the arms with the largest fractional parts, with ties going
to the earlier arm. The targets add up to $c$.

Before participant $i$, let $c$ be the smallest checkpoint with $c \ge i$, and
let $R = c - (i - 1)$ be the number of places left before it. The base
probability is
\begin{equation}
b_a(i) = \frac{\max\{t_a(c) - n_a(i), 0\}}{R},
\label{eq:base}
\end{equation}
which is the random allocation rule toward checkpoint $c$. The base
probabilities add up to one, and an arm whose target has been reached receives
probability zero. When only one arm is short of its target, $b_a(i) = 1$ for
that arm and the assignment is forced.

\subsection{The run-in}

The default run-in size is
\begin{equation}
N_1 = \min\{N, \; \lceil n_1 / s \rceil s\}, \qquad
n_1 = \max\{10, \lceil 0.10 N \rceil\},
\label{eq:runin}
\end{equation}
where $s = \sum_a \text{ratio}_a$ when every ratio is a whole number, and
$N_1 = \min\{N, n_1\}$ otherwise. Setting \texttt{runin(0)} removes the run-in.
For $i \le N_1$ the probability is $\phi(i) = \operatorname{tilt}(b(i), \hat a)$
of equation~(\ref{eq:tilt}), where the winner $\hat a$ comes from the votes of
the balance variables only. Because $b(i)$ becomes a unit vector once all but
one arm have reached their run-in targets, the run-in ends with exactly the
target counts in every realization, yet the covariates influence each earlier
assignment. The option \texttt{runinmethod(rar)} removes the tilt and gives the
plain random allocation rule that \citet{zhao2015msb} suggested.

\subsection{After the run-in}

For $i > N_1$ the procedure applies three steps in this order.
\begin{enumerate}
\item \emph{Quota.} If an arm is forced by equation~(\ref{eq:base}), assign to
it.
\item \emph{Big stick bound.} Let $B$ be the size bound (default 3). If any arm
has $D_a(i) \le -B$, only those arms can be chosen, with probabilities
proportional to $-D_a(i)$ \citep{soares1983}. Otherwise, arms with
$D_a(i) \ge B$ receive probability zero and the remaining base probabilities are
rescaled to add up to one.
\item \emph{MSB coin.} If the bound changed nothing,
$\phi(i) = \operatorname{tilt}(b(i), \hat a)$, where the winner now includes the
group-size vote.
\end{enumerate}
The package labels each assignment with its step: \texttt{runin},
\texttt{quota}, \texttt{bound}, \texttt{msb-biased} (a winner existed), or
\texttt{msb-random} (no winner).

\subsection{Fixed internal choices}

Three choices are fixed constants rather than options. We did not study
alternatives to them, so the simulations below describe the procedure with
these constants.
\begin{itemize}
\item \emph{The group-size vote starts at a shortfall of 1.} A shortfall below
1 cannot be corrected by one assignment, and under 1:1 allocation a deviation of
$\pm 0.5$ occurs at every odd count, so a lower limit would make group size vote
almost constantly.
\item \emph{The group-size vote has weight 1.} Because the balance-variable
weights add up to one, group size then outvotes any proper subset of the
variables and ties only when every variable votes against it.
\item \emph{The base probability comes from the places left before the next
checkpoint.} This is what makes the checkpoints exact, and it moves every arm
toward its target throughout enrollment instead of leaving the whole correction
to the end.
\end{itemize}

\section{The commands}
\label{sec:commands}

\subsection{Overview}

All randomization commands are subcommands of \gracedr, and all take
\texttt{study(}\emph{folder}\texttt{)}, the folder that holds the study files.
Table~\ref{tab:commands} lists them. The dispatcher \texttt{gracedr} runs
\texttt{gracedr\_}\emph{subcommand}, so each command is a separate ado-file.

\begin{table}[htbp]
\centering
\caption{Commands of the \gracedr\ package}
\label{tab:commands}
\begin{tabular}{ll}
\hline
Command & Purpose \\
\hline
\texttt{gracedr setup} & Create a study folder from a design \\
\texttt{gracedr codes} & Check and load a file of treatment codes \\
\texttt{gracedr randomize} & Randomize one participant and save the record \\
\texttt{gracedr status} & Group sizes, balance tables, and graphs \\
\texttt{gracedr override} & Change an assignment with a recorded reason \\
\texttt{gracedr audit} & List the append-only audit log \\
\texttt{gracedr replay} & Redraw a past assignment and compare it with the
record \\
\texttt{gracedr simulate} & Simulate trials under up to five allocation
methods \\
\texttt{encsave} & Save the data in memory to an encrypted zip file \\
\texttt{encuse} & Load the data from an encrypted zip file \\
\hline
\end{tabular}
\end{table}

\subsection{Defining a study}

\begin{stlog}
\gracedr\ setup, study(\emph{folder}) n(\#) arms(\emph{string})
    [categorical(\emph{string}) continuous(\emph{namelist}) weights(\emph{numlist})
    ratio(\emph{numlist}) pstar(\#) xi(\#) runin(\#) runinmethod(tilted|rar)
    forcedat(\emph{numlist}) sizebound(\#) codes(\emph{file}) seed(\#) encrypt]
\end{stlog}

\noindent Table~\ref{tab:options} lists the options. The balance variables are
named in \texttt{categorical()}, with their levels, and \texttt{continuous()}.
The command checks every option, creates the folder, and writes the design,
the treatment codes, the empty participant table, the random number state, and
the audit log. It refuses to run if the folder already holds a study.

\begin{table}[htbp]
\centering
\footnotesize
\caption{Options of \texttt{gracedr setup}}
\label{tab:options}
\begin{tabular}{lll}
\hline
Option & Default & Meaning \\
\hline
\texttt{n()} & (required) & Total planned enrollment \\
\texttt{arms()} & (required) & Arm labels, at least two \\
\texttt{categorical()} & none & Categorical balance variables and their
levels \\
\texttt{continuous()} & none & Continuous balance variables \\
\texttt{ratio()} & 1 for all & Allocation ratio, any positive numbers \\
\texttt{weights()} & equal & Positive weights of the balance variables \\
\texttt{pstar()} & 0.30 & Imbalance probability limit $p^\ast$ \\
\texttt{xi()} & 0.65 & Biased coin probability $\xi$ \\
\texttt{runin()} & default & Run-in size; 0 removes the run-in \\
\texttt{runinmethod()} & tilted & Tilted run-in or plain random allocation rule \\
\texttt{forcedat()} & none & Enrollment totals at which the target counts are
forced \\
\texttt{sizebound()} & 3 & Big stick bound $B$ \\
\texttt{codes()} & none & Treatment codes file with columns \texttt{code} and
\texttt{arm} \\
\texttt{seed()} & random & Seed of the random number generator \\
\texttt{encrypt} & off & Store every dataset as an encrypted zip file \\
\hline
\end{tabular}
\end{table}

\subsection{Treatment codes, the study folder, and the random number state}

A list of treatment codes is kept in a file with one row per code and the two
columns \texttt{code} and \texttt{arm}. A participant's treatment code and id are
what the user copies into the trial database. The package ships an example list
of 200 codes. \texttt{gracedr codes} checks that every arm is valid and that the
codes are unique.

A study folder holds five datasets: the design, the participants, the codes, the
audit log, and the random number state. A write goes first to a file with the
extension \texttt{.tmp}, the current file is then copied to a file with the
extension \texttt{.bak}, and the new file replaces the current one. Stata has no
file rename, so the replacement is a copy.

Randomization uses Stata's default generator. Each command saves the user's
state, installs the study's state, makes its draw, and restores the user's state,
so a study neither depends on nor disturbs the random numbers of an interactive
session. The state is a string of thousands of characters, which is stored in a
\texttt{strL} variable. Each participant row stores the state in force before
that participant's draw.

A draw takes one uniform random number from the study's state and uses it to
choose the arm from the allocation probabilities. When the study has treatment
codes, the draw takes a second uniform number from the same state and uses it to
choose among the unused codes of the chosen arm. One stored state therefore
reproduces both the arm and the treatment code. The new state is not stored
until the end: \texttt{gracedr randomize} writes the participant, code, and audit
files first and then writes the new state, and it refuses to write the state if
the state in the study folder is no longer the one that the draw started from. If
an earlier write fails, the state is unchanged, and the command lists the files that
were already written. \texttt{gracedr replay} installs the stored state, redoes the
draw from the data before it, and compares the arm and the treatment code with
the record.

\subsection{Worked example}

We create a study for 40 participants with two arms and the balance variables
sex, race, and age, using the example codes and a seed.

\begin{stlog}
. gracedr setup, study("studies/example") n(40) arms(A B)
>     categorical(sex: male female | race: white black other) continuous(age)
>     codes("data/example_codes.csv") seed(20260611)
gracedr: 200 treatment codes are valid.
  arm A: 100 codes
  arm B: 100 codes
gracedr: study created in studies/example
  40 participants planned, run-in 10, seed 20260611
\end{stlog}

\noindent The default run-in is 10 participants, from equation~(\ref{eq:runin}).
We then randomize the 40 made-up participants of the file
\texttt{example\_participants.dta} one at a time. Each call prints the id, the
arm, the treatment code, and the group sizes. The last one is

\begin{stlog}
. gracedr randomize, study("studies/example") id(P040) sex(female) race(other)
>     age(71.3)
Participant P040 randomized to arm A
Treatment code: 11394
Group sizes
Arm                        n     pct   Next target  Final target
A                         20    50.0            20            20
B                         20    50.0            20            20
\end{stlog}

\noindent The counts are exact at the final total of 40.
\texttt{gracedr status} prints the group sizes and the balance table and, with
\texttt{graphs}, draws the two graphs of figure~\ref{fig:status}.

\begin{stlog}
. gracedr status, study("studies/example") graphs
Study status: 40 participants randomized
Group sizes
Arm                        n     pct   Next target  Final target
A                         20    50.0            20            20
B                         20    50.0            20            20
Balance by arm
  arm1 = A
  arm2 = B

  +-----------------------------------------------------------------+
  | variable    level          arm1         arm2   nmissing     smd |
  |-----------------------------------------------------------------|
  |      sex     male    13 (65.0%)   14 (70.0%)          0   0.107 |
  |      sex   female     7 (35.0%)    6 (30.0%)          0   0.107 |
  |     race    white    10 (50.0%)   14 (70.0%)          0   0.417 |
  |     race    black     4 (20.0%)    3 (15.0%)          0   0.417 |
  |     race    other     6 (30.0%)    3 (15.0%)          0   0.417 |
  |      age            60.0 (13.0)   60.0 (8.9)          0   0.000 |
  +-----------------------------------------------------------------+
\end{stlog}

\begin{figure}[htbp]
\centering
\includegraphics[width=0.48\textwidth]{status-imbalance.pdf}
\hfill
\includegraphics[width=0.48\textwidth]{status-smd.pdf}
\caption{The two graphs of \texttt{gracedr status} for the 40 participants.
Left: deviation of each arm from its target count as enrollment proceeds, with
the checkpoints marked. Right: largest absolute standardized mean difference of
each balance variable.}
\label{fig:status}
\end{figure}

An override changes one assignment and records a reason. It voids the old
treatment code and assigns an unused code of the new arm. The original draw
stays in the audit log.

\begin{stlog}
. gracedr override, study("studies/example") id(P040) arm(B)
>     reason("Wrong race entered in the example.")
Participant P040 changed from arm A to arm B
Code 11394 voided, new code 71284
\end{stlog}

\noindent The last two rows of the audit log are then

\begin{stlog}
  |    action     id   old   new                               reason |
  |-------------------------------------------------------------------|
  | randomize   P040           A                                      |
  |  override   P040     A     B   Wrong race entered in the example. |
\end{stlog}

\noindent A replay of each of the 40 participants reproduces the recorded arm
and code, including P040, whose comparison uses the original arm and code from
the audit log.

\subsection{Simulation}

\begin{stlog}
\gracedr\ simulate, methods(\emph{string}) nsims(\#) seed(\#)
    [study(\emph{folder}) | \emph{design options}] [generator(iid|drift)
    probs(\emph{string}) meansd(\emph{string}) driftto(\emph{string}) phigh(\#)
    saving(\emph{file}) summary(\emph{file}) clear]
\end{stlog}

\noindent The methods are \texttt{mmsb}, \texttt{srs} (simple randomization),
\texttt{blocks} (permuted blocks of size twice the sum of the ratio, for
whole-number ratios only),
\texttt{minimization} (the Pocock and Simon rule, \citealp{pocock1975}), and
\texttt{msb} (unmodified MSB: no group-size vote, no quota, and no bound).
Virtual participants come from independent distributions for each balance
variable: category probabilities for a categorical variable, given in
\texttt{probs()}, and a normal distribution for a continuous one, given in
\texttt{meansd()}. With \texttt{generator(drift)} each value moves in a straight
line from its start value to its \texttt{driftto()} value over the enrollment.
All methods allocate the same stream of participants, with the same random
numbers, so the results for one method do not depend on which other methods are
run. Simulation $j$ has its own seed drawn from \texttt{seed()}, so the first
$j$ simulations are the same whatever \texttt{nsims()} is. Simulations run one at
a time in a loop. With \texttt{study()} the command reads only the design and
writes nothing to the folder.

The results dataset has one row per simulation, method, and metric. The metrics
are the largest absolute standardized mean difference and the standardized
difference of each variable, and the absolute difference between each arm's
count and its target, each at the end of the run-in, at each forced total, and
at the end. Over the whole trial, they are the percent of assignments with a
probability of 1, the longest run of forced assignments to one arm, and the
expected probability of a correct guess for a guesser who always picks the arm
furthest behind its target, computed from the stored assignment probabilities
\citep{blackwell1957}.

\section{Encrypted datasets}
\label{sec:encrypt}

\subsection{Motivation}

The allocation sequence of a trial must stay concealed from the people who
enroll participants, so the files of a live study should be unreadable without a
password. Stata has no encrypted file format. We decided on a utility that is
separate from the randomization commands and does not know about them, built on
Python because Stata can run Python code and Python has a maintained package for
AES-encrypted zip files, \texttt{pyzipper}. Users must have installed Python
and that Python package.  They may also need to set the \texttt{python\_exec}
system variable to ensure that our new commands are able to run Python on 
the user's machine.

\subsection{encsave and encuse}

\begin{stlog}
encsave using \emph{filename} [, replace session]

encuse using \emph{filename} [, clear session]
\end{stlog}

\noindent \texttt{encsave} saves the dataset in memory to a temporary file,
zips it with a password, and erases the temporary file, so that only the zip file
remains. The zip uses AES-256 \citep{fips197}. \texttt{encuse} unzips to a
temporary file, loads the data, and erases the temporary file. Both erase the
temporary file on every exit path, including errors. A clear message tells the user how to
install \texttt{pyzipper} and which Python Stata is using.

\subsection{Handling the password}

The commands read the password from the environment variable
\texttt{GRACEDR\_PASSPHRASE}. If it is not set, a masked pop-up prompt opens
through Python and asks for it, and asks twice when it saves a file. There is
no \texttt{password()} option. Because Stata echoes command lines into the log and into
the command history, a password typed in a command would be written to
both. The password is not stored in a file, and it goes to Python in memory, not
on a command line. We check in the certification scripts that the password does
not appear in the log, in \texttt{set trace} output, or in a displayed macro. A
wrong password stops the command with error 499 and changes nothing.

For commands that read a study repeatedly, the \texttt{session} option keeps
the password in the memory of the Python process, once for each study folder, so
that the pop-up opens once per study and not once per command. The password is
forgotten after a wrong password, a different study folder gets its own entry,
and the environment variable always wins when it is set. Calling \texttt{encsave}
or \texttt{encuse} directly, without the \texttt{session} option, never caches the
password.

\subsection{Encrypted studies}

The option \texttt{encrypt} of \texttt{gracedr setup} stores the design, the
participants, the codes, the audit log, and the random number state as encrypted
zip files, written by \texttt{encsave}. Every other \texttt{gracedr} command
treats a folder that holds any of those zip files as an encrypted study, so it
needs no option, and they all use the \texttt{session} option. A write goes to a
file with the extension \texttt{.tmp.zip}, the current file is kept as a file
with the extension \texttt{.bak.zip}, and the new file then replaces it. Each
write first checks that the password opens the current archive, so a wrong
password cannot overwrite a study. A folder that holds any of the zip files can
never receive a plain dataset, because the extension is decided by what is
already in the folder.

In the worked example of section~\ref{sec:commands}, creating the same study
with the option \texttt{encrypt} and randomizing the same 40 participants gives
a folder with the files \texttt{audit.zip}, \texttt{codes.zip},
\texttt{design.zip}, \texttt{participants.zip}, and \texttt{state.zip}, plus a
\texttt{.bak.zip} copy of each file that was rewritten. With the same seed, the
same participants, and the same override, its participant table equals that of
the plain study in every variable except the time stamp.

\subsection{Limits}

Encryption protects files at rest. It does not hide the allocation from anyone
who holds the password, so the password should be shared only with those who are
meant to be unblinded. Stata itself makes temporary copies of the dataset in its
temporary folder while a command runs, and these copies are plain. Stata erases
them when the command ends, but after a crash one can remain. Users who need to
keep the allocation sequence from the machine's other users should set Stata's
temporary folder to a protected location. The commands that write the results of
a simulation or a table to a file named by the user write plain datasets,
because the user chooses those files.

\section{Simulation study}
\label{sec:sim}

\subsection{Design}

Each simulated trial enrolls $N \in \{200, 500\}$ participants into two arms in a
1:1 ratio. The balance variables are sex, race (three levels), and age, with
equal weights. The default run-in has 20 participants when $N = 200$ and 50 when
$N = 500$. Two enrollment streams are used. In the \emph{iid} stream every
participant is drawn from the same population: the probability of male sex is
0.48, race is white, black, or other with probabilities 0.62, 0.23, and 0.15,
and age is normal with mean 60 and standard deviation 10. In the \emph{drift}
stream the population changes steadily over enrollment: the probability of male
sex rises from 0.20 to 0.80, the probability of black race rises from 0.08 to
0.48 with the share of other race fixed at 0.15, and mean age rises from 48 to 72
years with standard deviation 8. Drift is a stress test, because any rule that
concentrates one arm early or late converts the time trend into covariate
imbalance.

Five methods allocate the same participants, with the defaults of the package.
The modified procedure and unmodified MSB use $p^\ast = 0.30$ and $\xi = 0.65$,
and the modified procedure also uses $B = 3$ and the default run-in. Permuted
blocks use blocks of size four. Minimization is the Pocock and Simon rule
\citep{pocock1975}: for each candidate arm it adds the participant to that arm,
takes the range of the arm counts divided by the target proportions within the
participant's level of each variable, and sums these ranges over the variables
with the design weights. Age is cut into three groups at the tertiles of 1,000
generated participants. The arm with the smallest total is assigned with
probability 0.8 (tied arms share it equally), and the other arms share the
remaining probability in proportion to their target proportions.

The measures fall into three groups. \emph{Balance} is the largest absolute standardized
mean difference (SMD) \citep{austin2009} over the three balance variables at the
end of enrollment and at the end of the run-in. For a categorical variable the
SMD is the largest over its levels, with each level treated as a 0/1 indicator.
\emph{Group size} is the absolute difference between the final count of the
first arm and its target, and the percentage of trials in which the counts are
exact at the end of the run-in and at the end of enrollment.
\emph{Predictability} is measured by the percentage of assignments that are
forced, meaning that some arm has probability one; by the longest run of
consecutive forced assignments to the same arm; and by the expected probability
of a correct guess under the strategy of \citet{blackwell1957}. The
correct-guess probability is computed from the stored probabilities $\phi(i)$,
not from realized draws.

Every setting uses 1,000 simulated trials and the seed 20260611. A Monte Carlo
standard error is the standard deviation of the per-trial values divided by
$\sqrt{1000}$. All methods allocate the same participants within a trial, so
differences between methods are computed within trials, and their standard
errors reflect that pairing.

\subsection{Results}

Figure~\ref{fig:tradeoff} summarizes balance and predictability.
Table~\ref{tab:balance} shows covariate balance. The modified procedure had a
mean largest SMD at the end of enrollment of 0.162, 0.159, 0.097, and 0.093 in
the four settings. This is lower than for simple randomization (0.212, 0.211,
0.136, and 0.133) and for permuted blocks (0.214, 0.198, 0.133, and 0.125), and
higher than for unmodified MSB (0.150, 0.152, 0.091, and 0.091) and for
minimization (0.077, 0.073, 0.040, and 0.037). Table~\ref{tab:paired} gives the
paired differences. Relative to unmodified MSB, the modified procedure was higher
by 0.002 to 0.012 in the mean largest SMD.

\begin{table}[htbp]
\centering
\footnotesize
\caption{Covariate balance. Largest absolute SMD over the three balance
variables at the end of enrollment (mean, with Monte Carlo standard error in
parentheses, and 95th percentile) and at the end of the run-in (mean). 1,000
simulated trials per setting.}
\label{tab:balance}
\begin{tabular}{llrrr}
\hline
Setting & Method & End (SE) & End, 95th pct. & Run-in \\
\hline
$N=200$, iid & Simple & 0.212 (0.003) & 0.359 & 0.768 \\
 & Permuted blocks & 0.214 (0.003) & 0.368 & 0.709 \\
 & Minimization & 0.077 (0.001) & 0.151 & 0.494 \\
 & MSB & 0.150 (0.001) & 0.221 & 0.622 \\
 & Modified MSB & 0.162 (0.002) & 0.254 & 0.617 \\
$N=200$, drift & Simple & 0.211 (0.003) & 0.355 & 0.734 \\
 & Permuted blocks & 0.198 (0.003) & 0.338 & 0.691 \\
 & Minimization & 0.073 (0.001) & 0.147 & 0.549 \\
 & MSB & 0.152 (0.001) & 0.228 & 0.628 \\
 & Modified MSB & 0.159 (0.002) & 0.251 & 0.613 \\
$N=500$, iid & Simple & 0.136 (0.002) & 0.225 & 0.446 \\
 & Permuted blocks & 0.133 (0.002) & 0.224 & 0.448 \\
 & Minimization & 0.040 (0.001) & 0.091 & 0.234 \\
 & MSB & 0.091 (0.001) & 0.136 & 0.341 \\
 & Modified MSB & 0.097 (0.001) & 0.148 & 0.342 \\
$N=500$, drift & Simple & 0.133 (0.002) & 0.213 & 0.447 \\
 & Permuted blocks & 0.125 (0.002) & 0.213 & 0.448 \\
 & Minimization & 0.037 (0.001) & 0.078 & 0.285 \\
 & MSB & 0.091 (0.001) & 0.135 & 0.350 \\
 & Modified MSB & 0.093 (0.001) & 0.143 & 0.353 \\
\hline
\end{tabular}
\end{table}

Table~\ref{tab:sizes} shows group sizes. The modified procedure met the group
sizes at the end of the run-in and at the end of enrollment in every trial in all
four settings. Permuted blocks of size four met the final sizes in every trial
but the run-in sizes in only 68\% of trials when $N = 500$, because the run-in
of 50 is not a multiple of the block size. Unmodified MSB missed the final sizes
by a mean of 5.75, 6.14, 9.23, and 10.22 participants, with a maximum of 44, and
it was exact at the end in 2.4\% to 5.3\% of trials. Minimization missed the
final sizes by a mean of 0.62 to 0.66 participants and was exact in 43.7\% to
46.8\% of trials.

\begin{table}[htbp]
\centering
\footnotesize
\caption{Group sizes. Percentage of trials in which the arm counts equal their
targets at the end of the run-in and at the end of enrollment (exact, in
percent), and the absolute
difference between the final count of the first arm and its target (mean, with
the maximum in parentheses). Methods: simple randomization (Simple), permuted
blocks of size four (Blocks), minimization, unmodified MSB, and the modified
procedure.}
\label{tab:sizes}
\begin{tabular}{llrrr}
\hline
Setting & Method & Run-in (\%) & End (\%) & Final difference \\
\hline
$N=200$, iid & Simple & 18.6 & 6.9 & 5.50 (23) \\
 & Permuted blocks & 100.0 & 100.0 & 0.00 (0) \\
 & Minimization & 48.9 & 45.3 & 0.64 (4) \\
 & MSB & 17.2 & 4.4 & 5.75 (25) \\
 & Modified MSB & 100.0 & 100.0 & 0.00 (0) \\
$N=200$, drift & Simple & 18.6 & 6.9 & 5.50 (23) \\
 & Permuted blocks & 100.0 & 100.0 & 0.00 (0) \\
 & Minimization & 54.8 & 46.7 & 0.63 (4) \\
 & MSB & 17.3 & 5.3 & 6.14 (26) \\
 & Modified MSB & 100.0 & 100.0 & 0.00 (0) \\
$N=500$, iid & Simple & 11.0 & 3.0 & 8.72 (37) \\
 & Permuted blocks & 68.0 & 100.0 & 0.00 (0) \\
 & Minimization & 46.2 & 43.7 & 0.66 (4) \\
 & MSB & 11.1 & 3.2 & 9.23 (34) \\
 & Modified MSB & 100.0 & 100.0 & 0.00 (0) \\
$N=500$, drift & Simple & 11.0 & 3.0 & 8.72 (37) \\
 & Permuted blocks & 68.0 & 100.0 & 0.00 (0) \\
 & Minimization & 50.7 & 46.8 & 0.62 (3) \\
 & MSB & 10.7 & 2.4 & 10.22 (44) \\
 & Modified MSB & 100.0 & 100.0 & 0.00 (0) \\
\hline
\end{tabular}
\end{table}

Table~\ref{tab:predictability} shows predictability. The modified procedure
forced 2.4\% of assignments when $N = 200$ and 1.7\% when $N = 500$, with a
longest run of forced assignments to one arm that averaged 2.1 to 2.3 and
reached 7 and 10. Its expected correct-guess probability was 59.34\%, 59.31\%,
58.50\%, and 58.48\%. This is 8.5 to 9.3 percentage points above simple
randomization, 11.4 to 12.4 points below permuted blocks of size four
(70.79\% to 70.85\%), and 2.4 to 3.4 points below minimization (61.66\% to
62.12\%). Unmodified MSB had a correct-guess probability of 49.59\% to 49.93\%.

\begin{table}[htbp]
\centering
\footnotesize
\caption{Predictability. Percentage of assignments that are forced, the longest
run of consecutive forced assignments to the same arm (mean, with the maximum in
parentheses), and the expected probability of a correct guess in percent under
the strategy of \citet{blackwell1957} (mean, with Monte Carlo standard error in
parentheses).}
\label{tab:predictability}
\begin{tabular}{llrrr}
\hline
Setting & Method & Forced (\%) & Longest run & Guess (\%) \\
\hline
$N=200$, iid & Simple & 0.0 & 0.0 (0) & 50.00 (0.00) \\
 & Permuted blocks & 33.2 & 2.0 (2) & 70.79 (0.03) \\
 & Minimization & 0.0 & 0.0 (0) & 61.75 (0.05) \\
 & MSB & 0.0 & 0.0 (0) & 49.89 (0.03) \\
 & Modified MSB & 2.4 & 2.1 (7) & 59.34 (0.04) \\
$N=200$, drift & Simple & 0.0 & 0.0 (0) & 50.00 (0.00) \\
 & Permuted blocks & 33.2 & 2.0 (2) & 70.79 (0.03) \\
 & Minimization & 0.0 & 0.0 (0) & 62.12 (0.05) \\
 & MSB & 0.0 & 0.0 (0) & 49.59 (0.05) \\
 & Modified MSB & 2.4 & 2.1 (7) & 59.31 (0.05) \\
$N=500$, iid & Simple & 0.0 & 0.0 (0) & 50.00 (0.00) \\
 & Permuted blocks & 33.4 & 2.0 (2) & 70.85 (0.02) \\
 & Minimization & 0.0 & 0.0 (0) & 61.66 (0.03) \\
 & MSB & 0.0 & 0.0 (0) & 49.93 (0.02) \\
 & Modified MSB & 1.7 & 2.2 (10) & 58.50 (0.03) \\
$N=500$, drift & Simple & 0.0 & 0.0 (0) & 50.00 (0.00) \\
 & Permuted blocks & 33.4 & 2.0 (2) & 70.85 (0.02) \\
 & Minimization & 0.0 & 0.0 (0) & 61.86 (0.04) \\
 & MSB & 0.0 & 0.0 (0) & 49.67 (0.04) \\
 & Modified MSB & 1.7 & 2.3 (10) & 58.48 (0.03) \\
\hline
\end{tabular}
\end{table}

\begin{table}[htbp]
\centering
\footnotesize
\caption{Differences from the modified procedure within trials (other method
minus modified procedure), with 95\% intervals from the paired Monte Carlo
standard error. A positive difference in the largest SMD means the other method
balances worse. A positive difference in the correct-guess probability means the
other method is more predictable.}
\label{tab:paired}
\begin{tabular}{llrr}
\hline
Setting & Method & Largest SMD & Guess (\%) \\
\hline
$N=200$, iid & Simple & 0.0499 (0.0452, 0.0546) &
$-$9.3 ($-$9.4, $-$9.3) \\
 & Permuted blocks & 0.0523 (0.0467, 0.0578) & 11.4 (11.3, 11.5) \\
 & Minimization & $-$0.0846 ($-$0.0884, $-$0.0807) & 2.4 (2.3, 2.5) \\
 & MSB & $-$0.0118 ($-$0.0152, $-$0.0085) & $-$9.4 ($-$9.6, $-$9.3) \\
$N=200$, drift & Simple & 0.0523 (0.0476, 0.0571) &
$-$9.3 ($-$9.4, $-$9.2) \\
 & Permuted blocks & 0.0390 (0.0336, 0.0444) & 11.5 (11.4, 11.6) \\
 & Minimization & $-$0.0859 ($-$0.0898, $-$0.0820) & 2.8 (2.7, 2.9) \\
 & MSB & $-$0.0072 ($-$0.0107, $-$0.0037) & $-$9.7 ($-$9.8, $-$9.6) \\
$N=500$, iid & Simple & 0.0396 (0.0367, 0.0425) &
$-$8.5 ($-$8.6, $-$8.4) \\
 & Permuted blocks & 0.0362 (0.0327, 0.0396) & 12.4 (12.3, 12.4) \\
 & Minimization & $-$0.0561 ($-$0.0585, $-$0.0538) & 3.2 (3.1, 3.2) \\
 & MSB & $-$0.0057 ($-$0.0076, $-$0.0039) & $-$8.6 ($-$8.6, $-$8.5) \\
$N=500$, drift & Simple & 0.0401 (0.0373, 0.0430) &
$-$8.5 ($-$8.5, $-$8.4) \\
 & Permuted blocks & 0.0317 (0.0286, 0.0349) & 12.4 (12.3, 12.4) \\
 & Minimization & $-$0.0558 ($-$0.0581, $-$0.0535) & 3.4 (3.3, 3.5) \\
 & MSB & $-$0.0021 ($-$0.0040, $-$0.0002) & $-$8.8 ($-$8.9, $-$8.7) \\
\hline
\end{tabular}
\end{table}

\begin{figure}[htbp]
\centering
\includegraphics[width=\textwidth]{tradeoff.pdf}
\caption{Balance and predictability in the four settings. Horizontal axis:
expected correct-guess probability. Vertical axis: mean largest SMD at the end of
enrollment. The lower left is better. Standard errors are smaller than the
plotting symbols.}
\label{fig:tradeoff}
\end{figure}

\section{Computational details}

The results in this article were produced with Stata 19.5 on Windows and the
\gracedr\ package installed from the GitHub repository
\texttt{BiostatDoc/graceDR}. The folder
\texttt{SJarticle/}\allowbreak\texttt{replication} holds the do-files that
install the package, install \texttt{pyzipper}, and reproduce every table,
figure, and listing of this article, and the log of the run that made them. The
whole run takes about nine minutes. The package is installed with the command
\texttt{net install gracedr} and its \texttt{from()} option, which names the
\texttt{stata} folder of the repository at the release tag \texttt{stata-v1.0.2}. The
certification scripts of the package are in the folder \texttt{stata/cert}.

\section{Conclusions}
\label{sec:conclusions}

The simulations support a specific and limited claim. The modified procedure
delivers what MSB as published does not: exact group sizes at the end of the
run-in and at the end of enrollment in every simulated trial. It does so with
1.7\% to 2.4\% forced assignments and a correct-guess probability 8.5 to 9.3
percentage points above simple randomization and 11.4 to 12.4 points below
permuted blocks of size four. The price in covariate balance, relative to
unmodified MSB, is small, 0.002 to 0.012 in the mean largest SMD at the end, and
the procedure kept most of the improvement of MSB over simple randomization.

The simulations do not show that the modified procedure balances covariates
better than every alternative. Minimization had clearly smaller standardized
differences in all four settings (tables~\ref{tab:balance} and
\ref{tab:paired}), with an expected correct-guess probability only 2.4 to 3.4
percentage points above that of the modified procedure, and its final group
sizes were exact in fewer than half of the trials. Which trade-off is preferable
depends on the protocol. If exact group sizes at named enrollment totals are
required, the modified procedure and permuted blocks guarantee them and
minimization as simulated here does not. If group sizes may differ by a few
participants and balance on the controlled covariates is the main concern,
minimization performed best on balance in these simulations.

Several limitations restrict how far the results extend.
\begin{itemize}
\item The simulations use one design: two arms in a 1:1 ratio and three balance
variables with equal weights. The software handles more arms, unequal ratios,
unequal weights, and forced totals before the end of enrollment, but none of
these were simulated here.
\item Only two enrollment streams were used, and the drift was a steady change
over the whole enrollment. Earlier, abrupt, or larger changes were not examined.
\item The modified procedure was run with its defaults, and the three fixed
internal choices were not varied. The comparators were also run at one setting
each, with blocks of size four and a probability of 0.8 in minimization.
\item Balance was measured by the largest SMD at the end of enrollment and at
the end of the run-in. We did not examine balance during enrollment, the power
of a treatment comparison, or bias in the estimated treatment effect, which
\citet{lauzon2022msb} studied for MSB. Minimization controls the margins that
the SMD measures, which may favor it on this measure.
\item The Monte Carlo standard errors describe simulation error for these
designs and streams, not uncertainty about other designs.
\item The correct-guess probability is the expected value under one guessing
strategy \citep{blackwell1957}. An adversary with more information could do
better.
\end{itemize}

Because live randomization and simulation use the same allocation code, an
investigator can run \texttt{gracedr simulate} on the trial's own design and
expected enrollment pattern before the trial starts. We recommend doing so
instead of relying on the settings simulated here. The encrypted study protects
the files at rest, but who can read the allocation still depends on who holds
the password.

\bibliographystyle{sj}
\bibliography{references}

\begin{aboutauthors}
Grace Huiyi Ma is a student in the STEM Magnet Program at Dutch Fork High
School in Irmo, South Carolina.

James W. Hardin is a professor of biostatistics in the Department of
Epidemiology and Biostatistics at the University of South Carolina in Columbia,
South Carolina.
\end{aboutauthors}

\end{document}
